% Generated draft: requires mathematical and source review.
\documentclass[11pt]{article}
\usepackage[margin=25mm]{geometry}
\usepackage{fontspec}
\setmainfont{Latin Modern Roman}
\usepackage{amsmath,amssymb,mathtools,unicode-math}
\setmathfont{Latin Modern Math}
\usepackage{enumitem,needspace,xurl,hyperref}
\setlength{\parindent}{0pt}
\setlength{\parskip}{5pt}
\setlength{\emergencystretch}{3em}
\newcommand{\N}{\mathbb{N}}
\newcommand{\Z}{\mathbb{Z}}
\newcommand{\Q}{\mathbb{Q}}
\newcommand{\R}{\mathbb{R}}
\newcommand{\C}{\mathbb{C}}
\newcommand{\F}{\mathbb{F}}
\newcommand{\A}{\mathbb{A}}
\newcommand{\PP}{\mathbb{P}}
\newcommand{\E}{\mathbb{E}}
\newcommand{\eps}{\varepsilon}
\newcommand{\dd}{\,\mathrm{d}}
\begin{document}
\section*{UnsolvedMath ID: 10600022}
\textbf{Source problem code:} catalogue entry
\subsection{Definitions and mathematical statement}
A rack is a set \(X\) with an operation \(x\triangleright y\) such that each right translation is bijective and
\[
 (x\triangleright y)\triangleright z
 =(x\triangleright z)\triangleright(y\triangleright z).
\]
In the standard rack-space convention, the chain group is \(C_n^R(X)=\Z[X^n]\), with boundary
\[
 \partial(x_1,\ldots,x_n)=\sum_{i=1}^n(-1)^i
 \bigl[(x_1,\ldots,\widehat{x_i},\ldots,x_n)
 -(x_1\triangleright x_i,\ldots,x_{i-1}\triangleright x_i,
                   x_{i+1},\ldots,x_n)\bigr].
\]
Its homology is the integral homology of the usual rack space. The problem is to develop the rack-space homology program and, in particular, investigate the properties and associated homology of the birack space. A birack involves an invertible nondegenerate set-theoretic Yang--Baxter crossing operation; the exact birack axioms and corresponding space construction are not specified in this excerpt, so no particular birack chain model is chosen.

The source reports geometric interpretations: two-dimensional cycles are represented by consistently rack-coloured virtual link diagrams, and three-dimensional cycles by such diagrams with regions also coloured, referring to Greene's thesis. The required crossing/region colour rules and framing conventions are omitted. It says second homology of the dihedral rack is treated in that thesis, and that for prime \(p\), third homology has a factor \(\Z_p\). The usual dihedral rack has underlying set \(\Z/p\Z\) and operation \(x\triangleright y=2y-x\). The source neither defines factor nor states whether \(p\) must be odd; the prime \(2\) requires checking, since that operation becomes trivial there. No missing second-homology formula or direct-summand assertion is invented.
\subsection{Short English statement}
Investigate rack-space homology through coloured virtual-link diagrams, and develop the corresponding birack-space and homology theory, retaining the stated but incompletely specified dihedral-rack results as background.
\subsection{Sources}
\begin{itemize}
\item \emph{Unsolved Problems in Virtual Knot Theory and Combinatorial Knot Theory}, revised survey, problem-list item 22, ``The Rack Space.'' \url{https://arxiv.org/abs/1409.2823}. Recorded access 2026-07-29; extraction source-tex.
\item The statement credits rack spaces to Fenn, Rourke and Sanderson, names further homology work of Carter, Kamada and Saito, and refers to Greene's thesis for coloured-cycle interpretations and second dihedral-rack homology. All these citation entries and the reference for the \(\Z_p\) factor are missing from the extraction.
\item Local UnsolvedMath registry, \texttt{data/unsolvedmath/problems.json}, permanent ID \(10600022\), exact source code \texttt{catalogue entry}; supplied partially\_solved/PARTIAL-PROGRESS status. Its summary reports the birack question unresolved without a newer precise citation.
\end{itemize}
\end{document}
